\begin{textsample}{POS Dim 1 – ro – Score 59.00 – ro/ro\_em\_15.txt}  \label{ex:f1_pos_137}
1.6.2.
Trabalho pedagógico com as metodologias ativas no ensino \textbf{médio} A expansão da era \textbf{digital} e a \textbf{demanda} por recursos tecnológicos presentes na sociedade da informação, do conhecimento ou da aprendizagem propõem \textbf{inúmeros} desafios ao trabalho pedagógico na \textbf{contemporaneidade} ( ALVES E SOUSA, 2016 ).
Dentre os desafios, a nova cultura de aprendizagem ativa para atender o perfil de estudante revela a importância da preparação do professor para o uso e a atuação com as ferramentas tecnológicas em diferentes espaços de aprendizagem e nas diferentes áreas de conhecimento contempladas no currículo para Educação Básica.
Outro ponto importante \textbf{diz} respeito que ao cenário educacional, \textbf{devido} à falta de interesse do estudante, o distanciamento da sala de \textbf{aula} e os \textbf{métodos} tradicionais de ensino que \textbf{desencadeiam} importantes reflexões acerca da significação da aprendizagem e sobre a inovação dos processos educativos.
Concomitante, entende -se que é \textbf{crucial} a formação continuada dos professores pautada em pressupostos epistemológicos \textbf{inerentes} a uma prática educativa mediadora, \textbf{crítica} e \textbf{reflexiva} na qual o estudante é o centro da aprendizagem ( FIALHO E MACHADO, 2017 ), o que corrobora com a relevância e \textbf{aplicabilidade} das metodologias ativas no Ensino Médio.
Nesse contexto, o trabalho pedagógico pautado em metodologias ativas no Ensino Médio é reflexo de um processo de metamorfose existente na sociedade \textbf{contemporânea} sob o qual tem sido vislumbrada a necessidade \textbf{emergente} de repensar o ensino, os \textbf{métodos} e as estratégias utilizadas pelo professor em sua prática educativa.
Como pressupostos iniciais, constituem fatores que culminam para emergência na transformação do ensino na Educação Básica, as mudanças globais no estilo de vida, no modo de viver, conhecer, ser, pensar, agir e aprender do estudante, e consequentemente, no ato de educar e formar essa nova geração.
O estudante da \textbf{contemporaneidade} requer uma formação que contemple as exigências da era \textbf{digital}, da comunicação e do conhecimento, sendo assim, mostra -se como essencial propiciar experiências educativas que o envolva como \textbf{pesquisador} e investigador para resolver \textbf{problemas} concretos que ocorrem no cotidiano de suas vidas, na sua comunidade e na sociedade de forma geral.
A aprendizagem precisa ser significativa, desafiadora, problematizadora e instigante, a ponto de \textbf{mobilizar} o estudante e o grupo a buscar soluções possíveis para serem discutidas e concretizadas à luz de referenciais \textbf{teóricos} práticos ( MORAN, MASETTO E BEHRENS, 2013 p.83 ).
Em \textbf{vista} disso, o referencial curricular do Estado de Rondônia aponta o campo das metodologias ativas com sendo um \textbf{caminho} \textbf{metodológico} e estratégico \textbf{relevante} e fundamental para reflexão na ação e sobre a ação ( SCHON, 2000 ) no que \textbf{diz} respeito a re(criar) experiências educativas significativas para o estudante do Ensino Médio nas áreas, entre as áreas de conhecimentos e na parte flexível do currículo.
Ademais, defende a importância da construção de uma prática educativa fundamentada na Educação por \textbf{competência} que forneça bases para que o professor possa compreender as mudanças nas concepções do aprender, planejar, ensinar e comunicar na sociedade da informação, do conhecimento e da aprendizagem.
Atrelada a isso, a produção do saber nas áreas do conhecimento \textbf{demanda} pelo desenvolvimento de ações que conduzem o professor e o estudante na busca dos processos de investigação e pesquisa com \textbf{vistas} a abrir \textbf{caminhos} coletivos para produção de conhecimento de forma integrada por parte do \textbf{aluno} e professor.
Nesse âmbito “ é necessário que o estudante ultrapasse o papel passivo de escutar, ler, decorar e de ser repetidor fiel dos ensinamentos do professor, para se tornar criativo, \textbf{crítico}, \textbf{pesquisador} e atuante para produzir conhecimento. ” ( MORAN, MASETTO E BEHRENS, 2013, p. 77 ).
Cabe destacar ainda que, as novas ferramentas tecnológicas presentes no mundo \textbf{contemporâneo}, além de \textbf{exigir} a ressignificação da prática educativa em sala de \textbf{aula} e o uso de novos ambientes de aprendizagem, também perpassam pela \textbf{compreensão} da aprendizagem situada em bases \textbf{teóricas} e práticas.
Desse modo, é \textbf{relevante} por parte do professor, a adoção de uma concepção \textbf{crítica} e \textbf{reflexiva} da ação docente que considera as especificidades dos estudantes no que \textbf{diz} respeito às mudanças sociais, políticas, financeiras, tecnológicas, e a peculiaridade regional e local.
Referente à perspectiva de aprendizagem significativa, a utilização de metodologias ativas no Ensino Médio favorece a autonomia, a curiosidade, a tomada de decisão individual e coletiva do estudante, e ainda, o protagonismo para o desenvolvimento de atividades que \textbf{permeiam} práticas sociais diferenciadas em conformidade com os contextos em que o estudante está inserido ( MORAN, MASETTO E BEHRENS, 2013 ; FIALHO E MACHADO, 2017 ).
Além disso, do ponto de \textbf{vista} da formação integral do estudante, conforme as \textbf{competências} gerais da BNCC, o trabalho pedagógico com as metodologias ativas contribui de forma direta no desenvolvimento social, conhecimento, criatividade, \textbf{capacidade} \textbf{crítica} \textbf{reflexiva}, autoconfiança, comunicação, cooperação, trabalho de equipe e empatia do estudante.
Sob essa perspectiva, as situações e os espaços de aprendizagem devem ser impulsionados pela curiosidade, pelo interesse, pela crise, pela \textbf{problematização} e pela busca de soluções possíveis para aquele momento histórico com a visão de que não há respostas únicas, absolutas e inquestionáveis ( MORAN, MASETTO E BEHRENS, 2013 ).
Nessa \textbf{vertente}, é necessário, segundo Freire ( 1996 ), que o docente perceba o educando enquanto sujeito \textbf{sócio-histórico} e cultural do ato de conhecer e amplie suas possibilidades educativas respeitando a coerência entre o saber-fazer e o saber ser.
Em \textbf{suma}, é \textbf{relevante} a construção de ambientes favoráveis à produção de conhecimento e desenvolvimento da autonomia do estudante.
Sendo assim, experiências como atividades realizadas em grupos com mais de um professor de classe acompanhando a execução de \textbf{tarefas}, realização de projetos, soluções de \textbf{problemas} reais e estudo de caso são estratégias que, se bem conduzidas, podem gerar uma verdadeira inovação pedagógica.
Cabe evidenciar ainda, que o ato de inovar em espaços de aprendizagem e na própria sala de \textbf{aula} amplia as possibilidades para estabelecer relações significativas entre os diferentes saberes, de \textbf{maneira} progressiva e mais elaborada, bem como converte a escola em lugar mais democrático, atrativo e estimulante ( DARO, 2018 ).
Em \textbf{suma}, é essencial no trabalho pedagógico, a \textbf{mobilização} e articulação de saberes considerando as experiências reais e/ou simuladas do estudante ( SOUZA, VILAÇA E TEIXEIRA, 2020 ) no \textbf{intuito} de oportunizar as condições para que possa solucionar \textbf{problemas} e \textbf{demandas} \textbf{complexas} presentes no seu cotidiano.
Sendo assim, no campo da \textbf{aplicabilidade} e do trabalho com as metodologias ativas nos espaços de aprendizagem, considerando a diversidade de estratégias, Camargo ( 2018b ) propõe, sob a perspectiva da sala de \textbf{aula} inovadora, estratégias para \textbf{fomentar} o aprendizado ativo, dentre as quais se destacam : a aprendizagem \textbf{baseada} em \textbf{problemas}, em projetos, a aprendizagem entre pares ou times, design thinking, gamificação, dentre outros.
Em relação ao trabalho com Design Thinking, Cavalcanti e Filatro ( 2016 ) afirmam que enquanto estratégias de aprendizagem correspondem a “ um modo de pensar ” \textbf{métodos} e estratégias para produção criativa de soluções inovadoras e promoção da centralidade do estudante no processo educativo.
Para \textbf{tanto}, constituem elementos chaves para as necessidades e as experiências do estudante, desde que imbuídos em três momentos : a inspiração, a ideação e a implementação.
Nesse contexto, a \textbf{figura} 5 \textbf{ilustra} um organograma com as etapas para criação de uma atividade com base na estratégia do Design Thinking, seguindo a técnica HCD Toolkit que \textbf{diz} respeito a hear ( ouvir ) ; create ( criar ) ; deliver ( entregar ).
Mediante à \textbf{figura} 4, cabe ressaltar que o Design Thinking na condição de estratégia propicia ao estudante criar opções e fazer escolhas que contribuam na produção de novas ideias e na \textbf{projeção} de soluções inovadoras.
Outro ponto de destaque, \textbf{diz} respeito que por meio da prototipagem, é possível tornar as ideias tangíveis para que sejam testadas e que permitam a discussão e/ou \textbf{análise} dos pontos fortes e fraquezas no desafio proposto.
Outra aplicação prática das metodologias ativas proposta por Flora ( 2015 ) é a gamificação.
Para a autora, a gamificação na Educação é uma estratégia de games ( jogos ), utilizada para promover o engajamento no ambiente de aprendizagem que por sua funcionalidade, revela -se como sendo um recurso que não pode faltar na caixa de ferramentas do professor, uma vez que \textbf{auxilia} no alcance dos objetivos de aprendizagem de forma divertida e propõe como \textbf{foco} \textbf{central} o desenvolvimento de estratégias para resolução de \textbf{problemas} que contribuem efetivamente para a aprendizagem significativa do estudante.
Nesta \textbf{vertente}, para o planejamento do trabalho pedagógico com estratégias gamificadas o organograma abaixo ( \textbf{figura} 6 ) orienta com base em Flora ( 2015 ) aspectos \textbf{relevantes} que devem ser observados quando a meta ou propósito do game ( jogo ) é a aprendizagem. \textbf{Figura} 6.
Organograma explicativo dos aspectos \textbf{relevantes} para o planejamento de uma atividade gamificada para o estudante do Ensino Médio no Estado de Rondônia. | Aspecto | \textbf{Descrição} | |---|---| | Meta | Motiva a ação do estudante ; a execução de \textbf{tarefas} ; a solução de \textbf{problemas} ; as mudanças de comportamento ; e o grau de dificuldade. | | Engajamento | Compartilhar o conhecimento ; contribuir para alcançar resultado ; criar algo que seja interessante e envolvente. | | Resolução de \textbf{problemas} | Oportuniza pensar sobre um \textbf{problema} ou atividade do dia a dia e convertê -lo em uma atividade que contenha os elementos do jogo ( competição, cooperação, exploração e premiação ). | | Pensamento de jogo | \textbf{Desperta} o interesse e prazer em jogar ; e a percepção de que participa e interage na construção de algo enquanto interagem. | | Diversão | Intensifica por meio de games a natureza competitiva, competitiva, o jogo com pessoas diferentes, se aceita o feedback. | Fonte : Adaptado de Flora ( 2015 ) Perante o organograma apresentado na \textbf{figura} 6, percebe -se a necessidade de articulação dos elementos do game, o pensamento de jogos e técnicas de design de forma dinâmica para motivar a ação, interação e o feedback no game realizado.
Sob a perspectiva da inovação e \textbf{progressão} da aprendizagem, destaca -se a participação e o envolvimento do estudante em atividades gamificadas com o uso de QR Code, pois oportunizam ao mesmo decifrar pistas e missões contidas nesses códigos, ao mesmo tempo em que \textbf{auxilia} no desenvolvimento da criatividade, autonomia, diálogo, na resolução de \textbf{problemas} e pode ser usado em componentes curriculares distintos.
Face ao \textbf{exposto}, chama a atenção para o fato de que a estratégia de gamificação no campo da Educação não é aplicada apenas com o uso de \textbf{tecnologias}, pois existem formas mais primitivas e comuns ao ambiente escolar que o professor sequer imagina que possa ser atividade gamificada.
Tal afirmação encontra fundamentação quando se analisam os jogos populares, jogos de tabuleiros e jogos pedagógicos realizados em sala de \textbf{aula} para diversão, engajamento e alcance de metas.
Como exemplo disso, no planejamento de uma atividade gamificada recomenda -se : a criação de um roteiro a ser cumprido, juntamente com a as missões a serem \textbf{desvendadas} pelos estudantes e, após a experiência gamificada, o professor pode realizar discussões com apresentação de ideias no \textbf{quadro} ou esquemas em papéis, incentivando os estudantes a criar de \textbf{maneira} colaborativa.
Por fim, no \textbf{intuito} de ampliar as possibilidades educativas para o trabalho pedagógico do professor com os estudantes do Ensino do Estado de Rondônia optou -se por apresentar no \textbf{quadro} 2 uma proposição de \textbf{detalhada} com exemplos de estratégias de metodologias ativas para o trabalho pedagógico e a \textbf{mobilização} dos saberes dos estudantes de forma ativa e significativa.

% matched lemmas: aluno, análise, aplicabilidade, aula, auxiliar, basear, caminho, capacidade, central, competência, complexo, compreensão, contemporaneidade, contemporâneo, crucial, crítico, demanda, descrição, desencadear, despertar, desvendar, detalhar, devido, digital, dizer, emergente, exigir, exposto, figura, figurar, foco, fomentar, ilustrar, inerente, intuito, inúmero, maneira, metodológico, mobilizar, mobilização, médio, método, permear, pesquisador, problema, problematização, progressão, projeção, quadro, reflexivo, relevante, suma, sócio-histórico, tanto, tarefa, tecnologia, teórico, vertente, vista
\end{textsample}
